\documentclass[10pt,a4paper]{amsart}
\usepackage[margin=19mm]{geometry}
\usepackage{amsmath,amssymb,mathtools}
\usepackage[scr=rsfs,cal=euler]{mathalfa}
\usepackage{xfrac}
\usepackage[unicode,hidelinks,bookmarks=false]{hyperref}
\newcommand{\ie}{i.e.\ }
\newcommand{\cf}{cf.\ }
\newcommand{\resp}{respectively\ }
\newcommand{\Subsection}{\subsection}
\providecommand{\nobreakdash}{\nobreak\hbox{-}}
\setlength{\emergencystretch}{2em}
\multlinegap0pt
\usepackage{mathrsfs}
\usepackage{stmaryrd}
\allowdisplaybreaks%
 \newtheorem{theorem}{Theorem}
 \newtheorem{corollary}{Corollary}
 \newtheorem{lemma}{Lemma}% 
\newtheorem{remark}{Remark}
\newtheorem{definition}{Definition}
\newtheorem{example}{Example}
\begin{document}
\title[On two conjectures on triangulations of 2-manifolds]{On two conjectures on triangulations of 2-manifolds}
\author{John M. Campbell}
\date{}
\maketitle
\noindent\emph{Source and licence.} On two conjectures on triangulations of 2-manifolds. Version arXiv:2608.06863v1 (CC BY 4.0). This material is distributed under the Creative Commons Attribution 4.0 International licence, http://creativecommons.org/licenses/by/4.0/, as recorded in the arXiv OAI licence metadata for this exact version and shown on the abstract page https://arxiv.org/abs/2608.06863v1. Full rights notice: licenses/arxiv-2608.06863-CC-BY-4.0.txt.\par\smallskip
\noindent\textit{Editorial scope.} Counted unit: the original \texttt{theorem} environment \texttt{firstsolution} at source lines 197--212 together with its complete proof at lines 214--385; the delivered document keeps source lines 144--386 so that every referenced label and every citation resolves inside it. Native editable source is the paper's own concatenated TeX; the AMS wrapper is editorial formatting only. No publisher class, upstream script or unrelated artwork is redistributed.
\par\smallskip\noindent\textit{Reference note.} This article ships no compiled bibliography (biblatex); the entries delivered here are composed from the author's own BibTeX (.bib) fields, with no field invented and no entry dropped.
\begin{equation}\label{EulerPoincare}
 \chi(T) = |V(T)| - |E(T)| + |F(T)|. 
\end{equation}
 If $T$ is a triangulation of $M$ (a 2-manifold), 
 then we write $\chi(T) = \chi(M)$. This is referred to as the \emph{Euler--Poincar\'{e} characteristic} of $M$. 

 Again for a finite abstract simplicial complex $T = (V(T), \mathcal{S}(T))$ and for $v \in V(T)$, the \emph{link} of $v$ in $T$
 refers to the subcomplex
\begin{equation}\label{definelk} 
 \operatorname{lk}_{T}(v) := \{ \sigma \in T : v \not\in \sigma \ \text{and} \ \sigma \cup \{ v \} \in T \}. 
\end{equation}
 The \emph{1-skeleton} of $T$ is denoted with $T^{(1)}$
 and may be defined so that 
\begin{equation}\label{defineskeleton}
 T^{(1)} := \{ \sigma \in T : \dim \sigma \leq 1 \}. 
\end{equation}

\subsection{Geometric realizations}
 Again let $T = \big( V(T), \mathcal{S}(T) \big)$ be a finite abstract simplicial complex. 
 We regard $\mathbb{R}^{V(T)}$ as the real vector space consisting of functions $x\colon V(T) \to \mathbb{R}$.
 For $v \in V(T)$, we write $e_{v} \in \mathbb{R}^{V(T)}$
 in place of the indicator function such that $e_{v}(w) = 1$ if $w = v$ and such that 
 $e_{v}(w)$ vanishes otherwise. For $\sigma \in \mathcal{S}(T)$, its \emph{geometric realization} is 
 $$ \left| \sigma \right| 
 = \left\{ \sum_{v \in \sigma} \lambda_{v} e_{v} : \lambda_{v} \geq 0 \ \text{for every} \ v \in \sigma, 
 \sum_{v \in \sigma} \lambda_{v} = 1 \right\}. $$
 Similarly, the \emph{geometric realization} of $T$ is the topological space
 $$ |T| := \bigcup_{\substack{\sigma \in \mathcal{S}(T) \\ \sigma \neq \varnothing }} 
 \left| \sigma \right| \subseteq \mathbb{R}^{V(T)} $$
 endowed with the subspace topology
 inherited from $\mathbb{R}^{V(T)}$ as a Euclidean space. 

\subsection{Related notions}
 A \emph{simplicial circle} is a 1-dimensional and finite abstract simplicial complex $C$ such that the geometric
 realization of $C$ is homeomorphic to $S^1$. Equivalently, the simplicial complex $C$ is isomorphic to the 
 boundary complex of a polygon. 

 By again letting $M$ denote a connected, closed 2-manifold, a 
 \emph{triangulation} $T$ of $M$
 is understood, for the purposes of this paper, as a finite simplicial triangulation of $M$, i.e., 
 so that $T$
 is a finite simplicial complex such that the geometric realization of $T$ is homeomorphic to $M$. 
 It should be emphasized that generalized and ideal triangulations are not considered in this paper. 
 For a finite abstract simplicial complex $D = (V(D), \mathcal{S}(D))$, 
 this is said to be a \emph{simplicial $2$-disk} if $D$ is a pure $2$-dimensional simplicial complex whose geometric realization is 
 homeomorphic to the closed unit disk 
 $\mathbb{D}^{2} = \{ x \in \mathbb{R}^{2} : \| x \| \leq 1 \}$. 

\section{Proof of Chen and Lawrencenko's second conjecture}\label{secproof}
 Our proof of the second Chen--Lawrencenko conjecture \cite[Conjecture 2]{ChenLawrencenko1999} provides a 
 stronger result, in the sense that we show how the required constant $C(M)$ 
 can be chosen to depend only on the Euler--Poincar\'{e} characteristic $\chi(M)$. 


\begin{theorem}\label{firstsolution}
 For every connected, closed 2-manifold $M$, 
 there exists a constant $C(M)$ such that: For every finite simplicial triangulation $T$ of $M$, 
 the relation 
\begin{equation}\label{equivalentconj}
 \xi(T) \leq \left| V(T) \right| + C(M) 
\end{equation}
 holds. Explicitly, by letting 
 \begin{equation}\label{defineKM}
 K_{M} := 18 \big( 1 + |\chi(M)| \big), 
\end{equation}
 one may then choose $C(M)$ as 
\begin{equation}\label{defineCM} 
 C(M) := K_{M}^{2} + 3 K_{M}. 
\end{equation} 
\end{theorem}

\begin{proof}
 For $v \in V(T)$, we write $d(v)$ in place of the number of triangular faces incident with $v$. The link of $v$ (recalling \eqref{definelk}) is 
 necessarily a simplicial circle, and this follows from $T$ being a simplicial triangulation of a (connected) 
 closed 2-manifold. We thus find that $d(v)$ is 
 the number of neighbors of $v$ in $T^{(1)}$ (recalling \eqref{defineskeleton}), 
 i.e., so that 
\begin{equation}\label{crudedv} 
 d(v) \leq \left| V(T) \right| - 1. 
\end{equation} 

 Let $u$ and $v$ denote distinct vertices in $V(T)$. 
 If a face $\sigma \in F(T)$ contains each of $u$ and $v$, 
 then the closure property associated with simplicial complexes gives us that 
 $\{ u, v \} \in E(T)$. From the assumption that $T$ triangulates $M$, and since $M$ is a (connected) closed 2-manifold (and thus does 
 not have a boundary), 
 every edge is necessarily contained in precisely two faces, i.e., so that
 \begin{equation*}
 \left| \{ \sigma \in F(T) : \{ u, v \} \subseteq \sigma \} \right| = \begin{cases} 
 2, & \text{if $\{ u, v \} \in E(T)$;} \\ 
 0, & \text{otherwise. } 
 \end{cases} 
\end{equation*}
 Since each edge in $E(T)$ is incident with precisely two faces in $F(T)$, 
 the equality 
\begin{equation}\label{3FT2ET} 
 3 |F(T)| = 2 |E(T)| 
\end{equation}
 holds. Write $\chi = \chi(M)$. From \eqref{3FT2ET} together with the 
 formulation of the Euler--Poincar\'{e} formula in \eqref{EulerPoincare}, 
 we obtain that 
\begin{equation}\label{ET3VTchi} 
 \left| E(T) \right| = 3 \left( \left| V(T) \right| - \chi \right). 
\end{equation}
 By the handshaking lemma (in the usual graph-theoretic sense) 
 together with \eqref{3FT2ET} and \eqref{ET3VTchi}, we find that 
\begin{equation}\label{fromhandshake} 
 \sum_{v \in V(T)} d(v) = 3 \left| F(T) \right| = 6 \left( \left| V(T) \right| - \chi \right). 
\end{equation}
 Define 
\begin{equation}\label{defineHT} 
 H(T) := \left\{ v \in V(T) : d(v) > \frac{ \left| V(T) \right| }{3} \right\}. 
\end{equation}
 From \eqref{fromhandshake} together with the definition on display in \eqref{defineHT}, we find (being mindful for the case whereby
 $H(T) = \varnothing$) that 
\begin{equation}\label{boundhandshake} 
 \left| H(T) \right| \frac{ \left| V(T) \right| }{3} 
 \leq \sum_{v \in H(T)} d(v) \leq 6 \left( \left| V(T) \right| - \chi \right). 
\end{equation}
 In turn, the bounds in \eqref{boundhandshake} 
 give us that 
\begin{equation}\label{boundHT} 
 \left| H(T) \right| \leq 18 \left( 1 - \frac{\chi}{ \left| V(T) \right| } \right). 
\end{equation}
 Since $\left| V(T) \right| \geq 1$, we see that 
 $1 - \frac{\chi}{ \left| V(T) \right| } \leq 1 + \left| \chi \right|$. 
 So, by defining $K_{M}$ as in \eqref{defineKM}, 
 the relation \eqref{boundHT} gives us that 
\begin{equation}\label{HTleqKM} 
 \left| H(T) \right| \leq K_{M}.
\end{equation} 
 So, the number of vertices of degree greater than $\frac{\left| V(T) \right|}{3}$
 is bounded above by a constant depending only on $M$. 

 Define $C(M)$ as in \eqref{defineCM}. 
 Our goal, at this point, is to show that $\left| V(T) \right| + C(M)$
 colors suffice, i.e., in order for there to 
 be a cyclic coloration of $T$. 

 We write $\mathcal{E}$ in place of the set consisting of elements $f \in F(T)$ containing at least two vertices in $H(T)$. We then find that 
 there are at most $\binom{ \left| H(T) \right| }{2}$ 
 sets $\{ u, v \}$ of distinct vertices $u, v \in H(T)$
 such that $\{ u, v \}$ is contained in a face in $\mathcal{E}$, 
 and that an arbitrary 2-set $\{ u, v \}$ satisfying the given conditions
 is contained in at most two faces. Consequently, the relation 
\begin{equation}\label{double2sets} 
 \left| \mathcal{E} \right| \leq 2 \binom{ \left| H(T) \right| }{2} 
\end{equation}
 holds, so that \eqref{HTleqKM} and \eqref{double2sets} together give us that 
\begin{equation}\label{boundexceptional} 
 \left| \mathcal{E} \right| \leq K_{M} \left( K_{M} - 1 \right). 
\end{equation}
 We then assign $\left| \mathcal{E} \right|$
 distinct colors to the faces in $\mathcal{E}$. 

 Since we have colored the faces in $F(T)$ containing at least two vertices in $H(T)$, 
 this leads us to consider the faces in $F(T)$ containing exactly one vertex in $H(T)$.
 In this direction, for $x \in H(T)$, we define
\begin{equation}\label{definecalF}
 \mathcal{F}_{x} := \{ \sigma \in F(T) : \sigma \cap H(T) = \{ x \} \}. 
\end{equation}
 
 If $H(T) = \varnothing$, the following stage is vacuous; otherwise, we endow $H(T)$ with a fixed (but arbitrary) linear order relation 
 $\triangleleft$, letting the vertices of $H(T)$ be written as $x_{1}$, $x_{2}$, $\ldots$, $x_{|H(T)|}$
 and ordered so that 
\begin{equation}\label{displayorderH} 
 x_{1} \triangleleft x_{2} \triangleleft \ldots \triangleleft x_{|H(T)|}.
\end{equation}
 Our strategy, at this point, is to 
 color the faces in the families 
 $\mathcal{F}_{x_{1}}$, $\mathcal{F}_{x_{2}}$, $\ldots$, $\mathcal{F}_{x_{|H(T)|}}$
 according to the ordering in \eqref{displayorderH}, 
 so that all of the faces within $\mathcal{F}_{x_{i}}$ are colored in a fixed but arbitrary order
 for fixed $i \in \{ 1, 2, \ldots, |H(T)| \}$. 

 According to the above process, suppose that 
 the face 
\begin{equation}\label{definesigma} 
 \sigma = \{ x_{i}, a, b \} \in \mathcal{F}_{x_{i}}
\end{equation}
 is to be colored, for some index 
 $i \in \{ 1, 2, \ldots, |H(T)| \}$ 
 and some (distinct) vertices $a$ and $b$, with the understanding that all preceding faces have been colored.
 Observe that: By construction, the vertex $x_i$ is the unique vertex that is both in $\{ x_i, a, b \}$ and in $H(T)$. 
 Now, consider previously colored faces (prior to $\sigma$ being colored) 
  that are in the same family $\mathcal{F}_{x_{i}}$ as $\sigma$. 
 From the definition in \eqref{definecalF}, 
 we find that each face in $\mathcal{F}_{x_{i}}$ necessarily contains $x_{i}$. 
 So, the number of previously colored faces
 in $\mathcal{F}_{x_{i}}$
 is bounded above by 
\begin{equation}\label{boundprevious} 
 \left| \mathcal{F}_{x_{i}} \right| - 1 \leq d(x_i) - 1. 
\end{equation} 
 From \eqref{crudedv} and \eqref{boundprevious} together,  we then find that the number of previously colored faces in 
  $\mathcal{F}_{x_{i}}$ is bounded above by 
\begin{equation}\label{upperVTminus2}
 d(x_{i}) - 1 \leq \left| V(T) \right| - 2.
\end{equation}
 We also consider the possibility that previously colored faces may be in the exceptional set $\mathcal{E}$, recalling the upper bound on 
 display in \eqref{boundexceptional}. Moreover, we need to consider, as below, faces in previously colored families. 

 We remain consistent with our notation for $\sigma$, as defined in \eqref{definesigma}, with $i$ as in \eqref{definesigma}. Now, let
 $j < i$. Let $\tau \in \mathcal{F}_{x_{j}}$ be a face having a common vertex with $\sigma$. 
 Since $\tau \in \mathcal{F}_{x_j}$ we have $\tau \cap H(T) = \{ x_{j} \}$.
 In particular, we find that $x_i \not\in \tau$. 
 Hence, if $\tau$ intersects $\sigma = \{ x_i, a, b \}$, 
 then $\tau$ contains $a$ or $b$. 

 We then see that at most two faces in $F(T)$ contain both $a$ and $x_j$ and that at most two faces in $F(T)$ contain both $b$ and $x_j$. 
 Consequently, for each previously processed vertex $x_j$ in $H(T)$, at most four faces in $\mathcal{F}_{x_j}$ can have a common vertex 
 with $\sigma$. There are at most $|H(T)| - 1$ previously processed vertices in $H(T)$, 
 and \eqref{HTleqKM} then gives us that there are at most $K_M$ previously processed vertices in $H(T)$. 
 So, out of the families among $\mathcal{F}_{x_1}$, $\mathcal{F}_{x_2}$, $\ldots$, $\mathcal{F}_{x_{i-1}}$, 
 there are at most $4 K_M$ previously colored faces 
 sharing a common vertex with $\sigma$. 
 This, together with foregoing considerations (recalling \eqref{boundexceptional} and \eqref{upperVTminus2}), 
 allows us to conclude that the number of previously colored faces that cannot have the same color as $\sigma$
 is at most 
 $\left| V(T) \right| - 2 + K_{M} ( K_{M} - 1 ) + 4 K_M$. 
 From \eqref{defineCM}, this value is equal to 
 $\left| V(T) \right| - 2 + C(M)$. 

 So, according to the above procedure, at most $\left| V(T) \right| - 2 + C(M)$ colors are not permitted when $\sigma $ is being colored. 
 So, with a palette of $\left| V(T) \right| + C(M)$ colors, one may always color $\sigma$, without obstructions. Applying the greedy 
 process described above, 
 we may color every face containing exactly one vertex in $H(T)$. 
 
 It remains to consider elements in $F(T)$ containing no vertex in $H(T)$. Let these remaining faces be ordered in a fixed (but 
 arbitrary) way. For such a face $\sigma = \{ a, b, c \}$ currently being colored, 
 we find that: Since none of the vertices in $\{ a, b, c \}$ is in $H(T)$, 
 we have that each of $d(a)$, $d(b)$, and $d(c)$ is bounded above by $\frac{ \left| V(T) \right| }{3}$. 
 The number of faces other than $\sigma$ having a common vertex with $\sigma$ is at most 
 $\big( d(a) - 1 \big) + \big( d(b) - 1 \big) + \big( d(c) - 1 \big)$, 
 which, in turn, is at most $\left| V(T) \right| - 3$, from the given bounds on $d(a)$, $d(b)$, and $d(c)$. 
 So, regardless of the number of faces that have already been colored, 
 at most $\left| V(T) \right| - 3$ colors are forbidden while $\sigma$ is being colored. 
 Since this value is strictly less than $\left| V(T) \right| + C(M)$, there is at least one color available for $\sigma$. 
 
 So, after completing the above procedure, all of the faces have been colored, and, by our above construction, any two faces sharing a 
 common vertex are colored differently. So, our coloring is cyclic and uses at most $\left| V(T) \right| + C(M)$ colors, and hence the 
 upper bound in \eqref{equivalentconj}. 
\end{proof}


\begin{thebibliography}{99}
\bibitem[]{ChenLawrencenko1999} Chen, Beifang; Lawrencenko, Serge. A note on cyclic colorations and minimal triangulations. Yokohama Math. J. 47 (1999) 93--99
\end{thebibliography}
\end{document}
